% TOP_PROBLEM: {"id": 2029, "title": "Erdős Problem #282", "category_id": 1, "category": {"id": 1, "name": "number_theory", "display_name": "Number Theory", "description": "Properties of integers, prime numbers, Diophantine equations.", "slug": "number-theory", "order_index": 1, "created_at": "2026-07-31T15:26:25.670Z"}, "status": "open"}
% FIRST_POSTED: null
% ENTRY_KIND: statement_only
% Imported from: unsolvedmath_additional_500.tex; UnsolvedMath ID 2029
% !TeX program = lualatex
% Compile twice: lualatex 2029.tex
% Self-contained: no external bibliography, images, or \input files.
% Numeric section labels and visible numbers are original UnsolvedMath IDs.
\documentclass[11pt,a4paper]{article}
\usepackage[margin=24mm,headheight=14pt]{geometry}
\usepackage{fontspec}
\setmainfont{Latin Modern Roman}
\setsansfont{Latin Modern Sans}
\setmonofont{Latin Modern Mono}
\usepackage{amsmath,amssymb,mathtools}
\usepackage{unicode-math}
\setmathfont{Latin Modern Math}
\usepackage{microtype}
\usepackage{enumitem,xurl,needspace}
\usepackage[dvipsnames]{xcolor}
\usepackage{hyperref}
\hypersetup{unicode=true,colorlinks=true,linkcolor=MidnightBlue,urlcolor=MidnightBlue,
 pdftitle={UnsolvedMath Problem 2029},pdfauthor={Source-annotated compilation},bookmarksnumbered=true}
\usepackage{bookmark,tocloft,titlesec,fancyhdr}
\setlength{\cftsecnumwidth}{4.2em}
\cftsetpnumwidth{2.5em}
\cftsetrmarg{4em}
\titleformat{\section}{\Large\bfseries\raggedright}{\thesection}{0.7em}{}
\pagestyle{fancy}\fancyhf{}
\fancyhead[L]{\small\sffamily Additional UnsolvedMath statements}
\fancyhead[R]{\small Original catalogue IDs}
\fancyfoot[C]{\thepage}
\setlength{\parindent}{0pt}
\setlength{\parskip}{5pt plus 1pt}
\setlength{\emergencystretch}{3em}
\setcounter{tocdepth}{1}\setcounter{secnumdepth}{1}
\setlist[itemize]{leftmargin=3em,itemsep=4pt,topsep=4pt,parsep=0pt}
\allowdisplaybreaks
\newcommand{\N}{\mathbb{N}}
\newcommand{\Z}{\mathbb{Z}}
\newcommand{\Q}{\mathbb{Q}}
\newcommand{\R}{\mathbb{R}}
\newcommand{\C}{\mathbb{C}}
\newcommand{\F}{\mathbb{F}}
\newcommand{\A}{\mathbb{A}}
\newcommand{\PP}{\mathbb{P}}
\newcommand{\E}{\mathbb{E}}
\newcommand{\eps}{\varepsilon}
\newcommand{\dd}{\,\mathrm d}
\newcommand{\sref}[2]{\hyperlink{source:#1:#2}{[S#2]}}
\newif\ifshowcatalogue
\showcataloguetrue
\newcommand{\cataloguescope}[1]{\ifshowcatalogue
 \paragraph{Statement recorded in the supplied source.}
 #1\fi}
\begin{document}
% UnsolvedMath ID 2029; source catalogue code EP-282

\setcounter{section}{2028}

\section{Termination of Restricted Greedy Egyptian Fractions}\label{2029}

{\small\sffamily UnsolvedMath ID 2029 · Number Theory}

\subsection{Definitions and mathematical statement}

\paragraph{Shared notation.}Unless a section says otherwise, $\N=\{1,2,\ldots\}$,
$\Z$ is the ring of integers, $\Q,\R,\C$ are the rational, real, and complex
fields, $[N]=\{1,\ldots,\lfloor N\rfloor\}$, and $\log$ is the natural
logarithm. For real functions of a parameter tending to infinity,
$f\ll g$ and $f=O(g)$ mean $|f|\leq Cg$ eventually for some fixed $C$;
$f\gg g$ reverses this comparison; $f=o(g)$ means $f/g\to0$;
$f\sim g$ means $f/g\to1$; and $f\asymp g$ means both order bounds.
A subscript on $O$ or $\ll$ specifies allowed constant dependence.
The expression $n^{o(1)}$ in an upper bound means growth smaller than
$n^\epsilon$ for each fixed $\epsilon>0$. Local definitions govern any
exception, finite-field convention, or use of the same symbol for another
object. An asymptotic claim is not automatically an effective algorithm.



The odd-denominator condition refers to the rational number in lowest terms. At each stage the current positive remainder determines the next allowed denominator; reaching remainder zero terminates the algorithm. The set of permitted denominators is fixed throughout the expansion. \sref{2029}{1} \sref{2029}{2}

\paragraph{Statement recorded in the supplied source.}
Let $A\subseteq \mathbb{N}$ be an infinite set and consider the following greedy algorithm for a rational $x\in (0,1)$: choose the minimal $n\in A$ such that $n\geq 1/x$ and repeat with $x$ replaced by $x-\frac{1}{n}$. If this terminates after finitely many steps then this produces a representation of $x$ as the sum of distinct unit fractions with denominators from $A$.
Does this process always terminate if $x$ has odd denominator and $A$ is the set of odd numbers? More generally, for which pairs $x$ and $A$ does this process terminate? \sref{2029}{1}

\subsection{Short English statement}

When does the greedy reciprocal expansion terminate if only specified denominators are allowed, particularly odd denominators for rationals with odd reduced denominator?

\subsection{Sources}

{\small

\begin{itemize}

\item[\hypertarget{source:2029:1}{S1}] ULAM.AI, \emph{UnsolvedMath}, supplied \texttt{problems.json}, record 2029 (EP-282), statement and background. The embedded literature-review date is 2026-08-17; that is the archive’s review date, not a fresh certification. Dataset landing page: \url{https://huggingface.co/datasets/ulamai/UnsolvedMath}.

\item[\hypertarget{source:2029:2}{S2}] T. F. Bloom, \emph{Erdős Problems}, Problem 282, \url{https://www.erdosproblems.com/282}. Reference imported from the supplied record; the live problem page was not independently checked for this compilation.

\item[\hypertarget{source:2029:3}{S3}] R. L. Graham, On Finite Sums of Unit Fractions, Proceedings of the London Mathematical Society s3-14 (1964), 193-207, \nolinkurl{doi:10.1112/plms/s3-14.2.193.} \url{https://doi.org/10.1112/plms/s3-14.2.193}. Bibliographic information imported from the supplied record.

\item[\hypertarget{source:2029:4}{S4}] R. L. Graham, On Finite Sums of Reciprocals of Distinct nth Powers, Pacific Journal of Mathematics 14 (1964), 85-92. \url{https://msp.org/pjm/1964/14-1/pjm-v14-n1-p10-p.pdf}. Bibliographic information imported from the supplied record.

\item[\hypertarget{source:2029:5}{S5}] Graham, R. L., On finite sums of unit fractions. Proc. London Math. Soc. (3) (1964), 193-207. Bibliographic information imported from the supplied record.

\end{itemize}

}

\label{end:2029}
\end{document}
